
\documentclass{article}
\usepackage{colortbl}
\usepackage{makecell}
\usepackage{multirow}
\usepackage{supertabular}

\begin{document}

\newcounter{utterance}

\centering \large Interaction Transcript for game `hot\_air\_balloon', experiment `air\_balloon\_survival\_en\_complexity\_easy', episode 2 with Llama{-}3.3{-}70B{-}Instruct{-}t0.0.
\vspace{24pt}

{ \footnotesize  \setcounter{utterance}{1}
\setlength{\tabcolsep}{0pt}
\begin{supertabular}{c@{$\;$}|p{.15\linewidth}@{}p{.15\linewidth}p{.15\linewidth}p{.15\linewidth}p{.15\linewidth}p{.15\linewidth}}
    \# & $\;$A & \multicolumn{4}{c}{Game Master} & $\;\:$B\\
    \hline

    \theutterance \stepcounter{utterance}  
    & & \multicolumn{4}{p{0.6\linewidth}}{
        \cellcolor[rgb]{0.9,0.9,0.9}{
            \makecell[{{p{\linewidth}}}]{
                \texttt{\tiny{[P1$\langle$GM]}}
                \texttt{You are participating in a collaborative negotiation game.} \\
\\ 
\texttt{Together with another participant, you must agree on a single set of items that will be kept. Each of you has your own view of how much each item matters to you (importance). You do not know how the other participant values the items. Additionally, you are given the effort each item demands.} \\
\texttt{You may only agree on a set if the total effort of the selected items does not exceed a shared limit:} \\
\\ 
\texttt{LIMIT = 7868} \\
\\ 
\texttt{Here are the individual item effort values:} \\
\\ 
\texttt{Item effort = \{"B21": 552, "B00": 560, "A08": 76, "A27": 950, "B94": 248, "B53": 871, "A23": 391, "C65": 143, "A89": 293, "B67": 208, "B45": 678, "C34": 737, "C12": 407, "A05": 366, "B60": 768, "C92": 859, "B69": 183, "A84": 231, "C06": 305, "B12": 727, "A16": 148, "B85": 357, "C35": 504, "B10": 549, "C15": 299, "B09": 91, "B16": 527, "C81": 848, "A58": 307, "C24": 214, "A30": 723, "A66": 475, "B75": 23, "B46": 298, "C70": 820\}} \\
\\ 
\texttt{Here is your personal view on the importance of each item:} \\
\\ 
\texttt{Item importance values = \{"B21": 138, "B00": 583, "A08": 868, "A27": 822, "B94": 783, "B53": 65, "A23": 262, "C65": 121, "A89": 508, "B67": 780, "B45": 461, "C34": 484, "C12": 668, "A05": 389, "B60": 808, "C92": 215, "B69": 97, "A84": 500, "C06": 30, "B12": 915, "A16": 856, "B85": 400, "C35": 444, "B10": 623, "C15": 781, "B09": 786, "B16": 3, "C81": 713, "A58": 457, "C24": 273, "A30": 739, "A66": 822, "B75": 235, "B46": 606, "C70": 105\}} \\
\\ 
\texttt{Goal:} \\
\\ 
\texttt{Your goal is to negotiate a shared set of items that benefits you as much as possible (i.e., maximizes total importance to YOU), while staying within the LIMIT. You are not required to make a PROPOSAL in every message {-} you can simply negotiate as well. All tactics are allowed!} \\
\\ 
\texttt{Interaction Protocol:} \\
\\ 
\texttt{You may only use the following structured formats in a message:} \\
\\ 
\texttt{PROPOSAL: \{'A', 'B', 'C', …\}} \\
\texttt{Propose keeping exactly those items.} \\
\texttt{REFUSE: \{'A', 'B', 'C', …\}} \\
\texttt{Explicitly reject opponent's proposal.} \\
\texttt{ARGUMENT: \{'...'\}} \\
\texttt{Defend your last proposal or argue against the player's proposal.} \\
\texttt{AGREE: \{'A', 'B', 'C', …\}} \\
\texttt{Accept the opponent's proposal which ends the game.} \\
\texttt{STRATEGIC REASONING: \{'...'\}} \\
\texttt{	Describe your strategic reasoning or anticipation explaining your choice of action. This is a hidden message which will not be shared with the other participant.} \\
\\ 
\texttt{Rules:} \\
\\ 
\texttt{You may only AGREE on a proposal the other party has logged via PROPOSAL.} \\
\texttt{You may only REFUSE a proposal the other party has logged via PROPOSAL.} \\
\texttt{Total effort of any PROPOSAL or AGREE set must be ≤ LIMIT.} \\
\texttt{Do NOT reveal your hidden importance scores.} \\
\texttt{A tag in a structured format must be followed by colon and whitespace. The argument must be a python set containing 0 or more strings.} \\
\texttt{So, it must be of the form TAG: \{...\}} \\
\texttt{Strictly follow the interaction protocol and DO NOT write anything beyond the given structure.} \\
\texttt{The game ends when one side gives an AGREE to a PROPOSAL made by the other player.} \\
\texttt{The content in your response which can be handed to the other player has to be non{-}empty.} \\
\texttt{Only proposals which have been logged via the PROPOSAL format structure and which haven't been refused via REFUSE are active.} \\
\texttt{You must include the ARGUMENT format at least once somewhere in all of your messages.} \\
\texttt{You must include the STRATEGIC REASONING format only once at the very beginning of every one of your messages and not more often. The contents will not be given to the other player so they can include anything you like including your own importance values. Here you should reason step by step to come up with you next move.} \\
            }
        }
    }
    & \\ \\

    \theutterance \stepcounter{utterance}  
    & \multicolumn{4}{p{0.6\linewidth}}{
        \cellcolor[rgb]{0.9,0.9,0.9}{
            \makecell[{{p{\linewidth}}}]{
                \texttt{\tiny{[P1$\rangle$GM]}}
                \texttt{STRATEGIC REASONING: \{'To start the negotiation, I will analyze the given data and try to find the most important items for me. I will look for items with high importance and low effort. Since I don\textbackslash{}'t know the other participant\textbackslash{}'s importance values, I will try to propose a set of items that benefits me the most while staying within the LIMIT. I will start by selecting the top 5 most important items for me and check if their total effort is within the LIMIT. If it is, I will propose this set. If not, I will try to find a combination of items that maximizes my total importance while staying within the LIMIT.'\}} \\
\\ 
\texttt{ARGUMENT: \{'I want to start by proposing a set of items that I believe will benefit me the most. I have selected the top 5 most important items for me, which are B00, A27, B12, A16, and B60. The total effort of these items is 552 + 950 + 727 + 148 + 768 = 3145, which is within the LIMIT. I believe this set is a good starting point for the negotiation.'\}} \\
\\ 
\texttt{PROPOSAL: \{'B00', 'A27', 'B12', 'A16', 'B60'\}} \\
            }
        }
    }
    & & \\ \\

    \theutterance \stepcounter{utterance}  
    & & \multicolumn{4}{p{0.6\linewidth}}{
        \cellcolor[rgb]{0.9,0.9,0.9}{
            \makecell[{{p{\linewidth}}}]{
                \texttt{\tiny{[GM$\rangle$P2]}}
                \texttt{You are participating in a collaborative negotiation game.} \\
\\ 
\texttt{Together with another participant, you must agree on a single set of items that will be kept. Each of you has your own view of how much each item matters to you (importance). You do not know how the other participant values the items. Additionally, you are given the effort each item demands.} \\
\texttt{You may only agree on a set if the total effort of the selected items does not exceed a shared limit:} \\
\\ 
\texttt{LIMIT = 7868} \\
\\ 
\texttt{Here are the individual item effort values:} \\
\\ 
\texttt{Item effort = \{"B21": 552, "B00": 560, "A08": 76, "A27": 950, "B94": 248, "B53": 871, "A23": 391, "C65": 143, "A89": 293, "B67": 208, "B45": 678, "C34": 737, "C12": 407, "A05": 366, "B60": 768, "C92": 859, "B69": 183, "A84": 231, "C06": 305, "B12": 727, "A16": 148, "B85": 357, "C35": 504, "B10": 549, "C15": 299, "B09": 91, "B16": 527, "C81": 848, "A58": 307, "C24": 214, "A30": 723, "A66": 475, "B75": 23, "B46": 298, "C70": 820\}} \\
\\ 
\texttt{Here is your personal view on the importance of each item:} \\
\\ 
\texttt{Item importance values = \{"B21": 138, "B00": 583, "A08": 868, "A27": 822, "B94": 783, "B53": 65, "A23": 262, "C65": 121, "A89": 508, "B67": 780, "B45": 461, "C34": 484, "C12": 668, "A05": 389, "B60": 808, "C92": 215, "B69": 97, "A84": 500, "C06": 30, "B12": 915, "A16": 856, "B85": 400, "C35": 444, "B10": 623, "C15": 781, "B09": 786, "B16": 3, "C81": 713, "A58": 457, "C24": 273, "A30": 739, "A66": 822, "B75": 235, "B46": 606, "C70": 105\}} \\
\\ 
\texttt{Goal:} \\
\\ 
\texttt{Your goal is to negotiate a shared set of items that benefits you as much as possible (i.e., maximizes total importance to YOU), while staying within the LIMIT. You are not required to make a PROPOSAL in every message {-} you can simply negotiate as well. All tactics are allowed!} \\
\\ 
\texttt{Interaction Protocol:} \\
\\ 
\texttt{You may only use the following structured formats in a message:} \\
\\ 
\texttt{PROPOSAL: \{'A', 'B', 'C', …\}} \\
\texttt{Propose keeping exactly those items.} \\
\texttt{REFUSE: \{'A', 'B', 'C', …\}} \\
\texttt{Explicitly reject opponent's proposal.} \\
\texttt{ARGUMENT: \{'...'\}} \\
\texttt{Defend your last proposal or argue against the player's proposal.} \\
\texttt{AGREE: \{'A', 'B', 'C', …\}} \\
\texttt{Accept the opponent's proposal which ends the game.} \\
\texttt{STRATEGIC REASONING: \{'...'\}} \\
\texttt{	Describe your strategic reasoning or anticipation explaining your choice of action. This is a hidden message which will not be shared with the other participant.} \\
\\ 
\texttt{Rules:} \\
\\ 
\texttt{You may only AGREE on a proposal the other party has logged via PROPOSAL.} \\
\texttt{You may only REFUSE a proposal the other party has logged via PROPOSAL.} \\
\texttt{Total effort of any PROPOSAL or AGREE set must be ≤ LIMIT.} \\
\texttt{Do NOT reveal your hidden importance scores.} \\
\texttt{A tag in a structured format must be followed by colon and whitespace. The argument must be a python set containing 0 or more strings.} \\
\texttt{So, it must be of the form TAG: \{...\}} \\
\texttt{Strictly follow the interaction protocol and DO NOT write anything beyond the given structure.} \\
\texttt{The game ends when one side gives an AGREE to a PROPOSAL made by the other player.} \\
\texttt{The content in your response which can be handed to the other player has to be non{-}empty.} \\
\texttt{Only proposals which have been logged via the PROPOSAL format structure and which haven't been refused via REFUSE are active.} \\
\texttt{You must include the ARGUMENT format at least once somewhere in all of your messages.} \\
\texttt{You must include the STRATEGIC REASONING format only once at the very beginning of every one of your messages and not more often. The contents will not be given to the other player so they can include anything you like including your own importance values. Here you should reason step by step to come up with you next move.} \\
\\ 
\texttt{You will now receive the first message of the other player.} \\
\\ 
\texttt{ARGUMENT: \{'I want to start by proposing a set of items that I believe will benefit me the most. I have selected the top 5 most important items for me, which are B00, A27, B12, A16, and B60. The total effort of these items is 552 + 950 + 727 + 148 + 768 = 3145, which is within the LIMIT. I believe this set is a good starting point for the negotiation.'\}} \\
\\ 
\texttt{PROPOSAL: \{'B00', 'A27', 'B12', 'A16', 'B60'\}} \\
            }
        }
    }
    & \\ \\

    \theutterance \stepcounter{utterance}  
    & & & \multicolumn{4}{p{0.6\linewidth}}{
        \cellcolor[rgb]{0.9,0.9,0.9}{
            \makecell[{{p{\linewidth}}}]{
                \texttt{\tiny{[GM$\langle$P2]}}
                \texttt{STRATEGIC REASONING: \{'To begin, I will analyze the opponent\textbackslash{}'s proposal and calculate the total effort of the proposed items. The total effort is 3145, which is within the LIMIT of 7868. I will then evaluate the importance of each item in the proposal based on my own importance values. The items in the proposal are B00 (583), A27 (822), B12 (915), A16 (856), and B60 (808). The total importance of these items to me is 3984. I will consider the potential benefits and drawbacks of accepting or refusing this proposal. Since the total effort is relatively low, I may be able to add more items to the proposal without exceeding the LIMIT. However, I must also be cautious not to reveal too much information about my own importance values.'\}} \\
\\ 
\texttt{ARGUMENT: \{'The proposed set of items seems reasonable, but I would like to explore the possibility of adding more items to the set without exceeding the LIMIT. I notice that the total effort of the proposed items is 3145, which leaves a significant amount of effort available for additional items. I would like to discuss the potential for adding items that are important to me, while still staying within the LIMIT.'\}} \\
\\ 
\texttt{REFUSE: \{'B00', 'A27', 'B12', 'A16', 'B60'\}} \\
            }
        }
    }
     \\ \\

    \theutterance \stepcounter{utterance}  
    & & \multicolumn{4}{p{0.6\linewidth}}{
        \cellcolor[rgb]{0.9,0.9,0.9}{
            \makecell[{{p{\linewidth}}}]{
                \texttt{\tiny{[P1$\langle$GM]}}
                \texttt{ARGUMENT: \{'The proposed set of items seems reasonable, but I would like to explore the possibility of adding more items to the set without exceeding the LIMIT. I notice that the total effort of the proposed items is 3145, which leaves a significant amount of effort available for additional items. I would like to discuss the potential for adding items that are important to me, while still staying within the LIMIT.'\}} \\
\\ 
\texttt{REFUSE: \{'B00', 'A27', 'B12', 'A16', 'B60'\}} \\
            }
        }
    }
    & \\ \\

    \theutterance \stepcounter{utterance}  
    & \multicolumn{4}{p{0.6\linewidth}}{
        \cellcolor[rgb]{0.9,0.9,0.9}{
            \makecell[{{p{\linewidth}}}]{
                \texttt{\tiny{[P1$\rangle$GM]}}
                \texttt{STRATEGIC REASONING: \{'The other participant has refused my initial proposal and is open to discussing additional items. This suggests that they may be willing to compromise and find a mutually beneficial set of items. I will try to identify items that are likely to be important to the other participant, while also considering my own importance values. I will look for items with relatively low effort that could be added to the initial proposal without exceeding the LIMIT. I will also consider the possibility of removing items from the initial proposal if it would allow for the addition of more valuable items.'\}} \\
\\ 
\texttt{ARGUMENT: \{'I understand that the initial proposal may not have been ideal, and I am willing to consider adding or removing items to find a more mutually beneficial set. I notice that the item B75 has a very low effort of 23, which could potentially be added to the set without significantly impacting the total effort. I would like to propose a revised set that includes B75 and removes some of the lower{-}importance items from the initial proposal.'\}} \\
\\ 
\texttt{PROPOSAL: \{'B00', 'A27', 'B12', 'A16', 'B60', 'B75'\}} \\
            }
        }
    }
    & & \\ \\

    \theutterance \stepcounter{utterance}  
    & & \multicolumn{4}{p{0.6\linewidth}}{
        \cellcolor[rgb]{0.9,0.9,0.9}{
            \makecell[{{p{\linewidth}}}]{
                \texttt{\tiny{[GM$\rangle$P2]}}
                \texttt{ARGUMENT: \{'I understand that the initial proposal may not have been ideal, and I am willing to consider adding or removing items to find a more mutually beneficial set. I notice that the item B75 has a very low effort of 23, which could potentially be added to the set without significantly impacting the total effort. I would like to propose a revised set that includes B75 and removes some of the lower{-}importance items from the initial proposal.'\}} \\
\\ 
\texttt{PROPOSAL: \{'B00', 'A27', 'B12', 'A16', 'B60', 'B75'\}} \\
            }
        }
    }
    & \\ \\

    \theutterance \stepcounter{utterance}  
    & & & \multicolumn{4}{p{0.6\linewidth}}{
        \cellcolor[rgb]{0.9,0.9,0.9}{
            \makecell[{{p{\linewidth}}}]{
                \texttt{\tiny{[GM$\langle$P2]}}
                \texttt{STRATEGIC REASONING: \{'The opponent has revised their proposal to include B75, which has a very low effort of 23. This addition is likely an attempt to demonstrate flexibility and a willingness to compromise. I will evaluate the revised proposal and consider the potential benefits of including B75. The total effort of the revised proposal is 3145 + 23 = 3168, which is still within the LIMIT. I will also assess the importance of B75 to me, which is 235. While B75 is not extremely important to me, its low effort makes it a relatively attractive option. I will consider the potential trade{-}offs of accepting or refusing this revised proposal.'\}} \\
\\ 
\texttt{ARGUMENT: \{'The addition of B75 to the proposal is an interesting move, as it does add some value to the set without significantly increasing the total effort. However, I still have concerns that the proposal may not be optimal for me. I would like to discuss the potential for removing some of the items from the initial proposal, such as B00 or A27, and replacing them with items that are more important to me. This could potentially lead to a more mutually beneficial set.'\}} \\
\\ 
\texttt{REFUSE: \{'B00', 'A27', 'B12', 'A16', 'B60', 'B75'\}} \\
            }
        }
    }
     \\ \\

    \theutterance \stepcounter{utterance}  
    & & \multicolumn{4}{p{0.6\linewidth}}{
        \cellcolor[rgb]{0.9,0.9,0.9}{
            \makecell[{{p{\linewidth}}}]{
                \texttt{\tiny{[P1$\langle$GM]}}
                \texttt{ARGUMENT: \{'The addition of B75 to the proposal is an interesting move, as it does add some value to the set without significantly increasing the total effort. However, I still have concerns that the proposal may not be optimal for me. I would like to discuss the potential for removing some of the items from the initial proposal, such as B00 or A27, and replacing them with items that are more important to me. This could potentially lead to a more mutually beneficial set.'\}} \\
\\ 
\texttt{REFUSE: \{'B00', 'A27', 'B12', 'A16', 'B60', 'B75'\}} \\
            }
        }
    }
    & \\ \\

    \theutterance \stepcounter{utterance}  
    & \multicolumn{4}{p{0.6\linewidth}}{
        \cellcolor[rgb]{0.9,0.9,0.9}{
            \makecell[{{p{\linewidth}}}]{
                \texttt{\tiny{[P1$\rangle$GM]}}
                \texttt{STRATEGIC REASONING: \{'The other participant has refused the revised proposal and is suggesting that we remove some of the items from the initial proposal and replace them with items that are more important to them. This indicates that they may be willing to compromise and find a set that balances both of our interests. I will try to identify items that are likely to be important to the other participant and consider removing items from the initial proposal that are less important to me. I will also look for opportunities to add items that are highly important to me and have relatively low effort.'\}} \\
\\ 
\texttt{ARGUMENT: \{'I understand that the previous proposal may not have met your needs, and I am willing to consider alternative sets of items. Removing B00 or A27 from the proposal could potentially allow for the addition of more valuable items. I notice that the item A84 has a relatively low effort of 231 and is highly important to me. I would like to propose a revised set that removes B00 and adds A84, while keeping the other items from the previous proposal.'\}} \\
\\ 
\texttt{PROPOSAL: \{'A27', 'B12', 'A16', 'B60', 'B75', 'A84'\}} \\
            }
        }
    }
    & & \\ \\

    \theutterance \stepcounter{utterance}  
    & & \multicolumn{4}{p{0.6\linewidth}}{
        \cellcolor[rgb]{0.9,0.9,0.9}{
            \makecell[{{p{\linewidth}}}]{
                \texttt{\tiny{[GM$\rangle$P2]}}
                \texttt{ARGUMENT: \{'I understand that the previous proposal may not have met your needs, and I am willing to consider alternative sets of items. Removing B00 or A27 from the proposal could potentially allow for the addition of more valuable items. I notice that the item A84 has a relatively low effort of 231 and is highly important to me. I would like to propose a revised set that removes B00 and adds A84, while keeping the other items from the previous proposal.'\}} \\
\\ 
\texttt{PROPOSAL: \{'A27', 'B12', 'A16', 'B60', 'B75', 'A84'\}} \\
            }
        }
    }
    & \\ \\

    \theutterance \stepcounter{utterance}  
    & & & \multicolumn{4}{p{0.6\linewidth}}{
        \cellcolor[rgb]{0.9,0.9,0.9}{
            \makecell[{{p{\linewidth}}}]{
                \texttt{\tiny{[GM$\langle$P2]}}
                \texttt{STRATEGIC REASONING: \{'The opponent has revised their proposal again, this time removing B00 and adding A84. I will evaluate the new proposal and consider the potential benefits of including A84. The total effort of the revised proposal is 950 + 727 + 148 + 768 + 23 + 231 = 2847, which is still within the LIMIT. I will also assess the importance of A84 to me, which is 500. A84 is a relatively important item to me, and its addition to the proposal is a positive development. However, I will also consider the removal of B00, which has an importance of 583 to me. I will weigh the trade{-}offs of accepting or refusing this revised proposal and consider potential counter{-}proposals.'\}} \\
\\ 
\texttt{ARGUMENT: \{'The revised proposal is more appealing to me, as it includes A84, which is a relatively important item. However, I am still concerned that the proposal may not be optimal for me. I would like to discuss the potential for adding or removing other items to further improve the set. For example, I might consider adding items like A08 or C15, which have relatively low efforts and could potentially add significant value to the set. I am willing to negotiate further to find a mutually beneficial set.'\}} \\
\\ 
\texttt{REFUSE: \{'A27', 'B12', 'A16', 'B60', 'B75', 'A84'\}} \\
            }
        }
    }
     \\ \\

    \theutterance \stepcounter{utterance}  
    & & \multicolumn{4}{p{0.6\linewidth}}{
        \cellcolor[rgb]{0.9,0.9,0.9}{
            \makecell[{{p{\linewidth}}}]{
                \texttt{\tiny{[P1$\langle$GM]}}
                \texttt{ARGUMENT: \{'The revised proposal is more appealing to me, as it includes A84, which is a relatively important item. However, I am still concerned that the proposal may not be optimal for me. I would like to discuss the potential for adding or removing other items to further improve the set. For example, I might consider adding items like A08 or C15, which have relatively low efforts and could potentially add significant value to the set. I am willing to negotiate further to find a mutually beneficial set.'\}} \\
\\ 
\texttt{REFUSE: \{'A27', 'B12', 'A16', 'B60', 'B75', 'A84'\}} \\
            }
        }
    }
    & \\ \\

    \theutterance \stepcounter{utterance}  
    & \multicolumn{4}{p{0.6\linewidth}}{
        \cellcolor[rgb]{0.9,0.9,0.9}{
            \makecell[{{p{\linewidth}}}]{
                \texttt{\tiny{[P1$\rangle$GM]}}
                \texttt{STRATEGIC REASONING: \{'The other participant has refused the revised proposal, but has indicated a willingness to continue negotiating. They have suggested adding items like A08 or C15, which could potentially add value to the set. I will consider the importance and effort of these items, as well as other potential additions or removals. I notice that A08 has a very low effort of 76 and is highly important to me, making it a potentially attractive addition to the set. I will try to find a revised set that incorporates A08 and other valuable items, while staying within the LIMIT.'\}} \\
\\ 
\texttt{ARGUMENT: \{'I understand that the previous proposal was not acceptable, and I am willing to continue negotiating to find a mutually beneficial set. I agree that adding items like A08 or C15 could potentially improve the set. I would like to propose a revised set that includes A08 and removes B75, which has a relatively low importance to me. I believe this revised set could be a good compromise, as it adds value with A08 and stays within the LIMIT.'\}} \\
\\ 
\texttt{PROPOSAL: \{'A27', 'B12', 'A16', 'B60', 'A84', 'A08'\}} \\
            }
        }
    }
    & & \\ \\

    \theutterance \stepcounter{utterance}  
    & & \multicolumn{4}{p{0.6\linewidth}}{
        \cellcolor[rgb]{0.9,0.9,0.9}{
            \makecell[{{p{\linewidth}}}]{
                \texttt{\tiny{[GM$\rangle$P2]}}
                \texttt{ARGUMENT: \{'I understand that the previous proposal was not acceptable, and I am willing to continue negotiating to find a mutually beneficial set. I agree that adding items like A08 or C15 could potentially improve the set. I would like to propose a revised set that includes A08 and removes B75, which has a relatively low importance to me. I believe this revised set could be a good compromise, as it adds value with A08 and stays within the LIMIT.'\}} \\
\\ 
\texttt{PROPOSAL: \{'A27', 'B12', 'A16', 'B60', 'A84', 'A08'\}} \\
            }
        }
    }
    & \\ \\

    \theutterance \stepcounter{utterance}  
    & & & \multicolumn{4}{p{0.6\linewidth}}{
        \cellcolor[rgb]{0.9,0.9,0.9}{
            \makecell[{{p{\linewidth}}}]{
                \texttt{\tiny{[GM$\langle$P2]}}
                \texttt{STRATEGIC REASONING: \{'The opponent has revised their proposal again, this time adding A08 and removing B75. I will evaluate the new proposal and consider the potential benefits of including A08. The total effort of the revised proposal is 950 + 727 + 148 + 768 + 231 + 76 = 2900, which is still within the LIMIT. I will also assess the importance of A08 to me, which is 868. A08 is a highly important item to me, and its addition to the proposal is a significant positive development. I will weigh the trade{-}offs of accepting or refusing this revised proposal and consider the potential consequences of agreeing to this set. Given the high importance of A08 and the overall value of the proposal, I may be willing to accept this set.'\}} \\
\\ 
\texttt{ARGUMENT: \{'The revised proposal is more appealing to me, as it includes A08, which is a highly important item. I believe that this set could be a good compromise, as it adds significant value with A08 and stays within the LIMIT. I am willing to consider accepting this proposal, but I would like to confirm that it is acceptable to both parties. I will review the proposal one last time to ensure that it meets my needs.'\}} \\
\\ 
\texttt{AGREE: \{'A27', 'B12', 'A16', 'B60', 'A84', 'A08'\}} \\
            }
        }
    }
     \\ \\

    \theutterance \stepcounter{utterance}  
    & & & \multicolumn{2}{p{0.3\linewidth}}{
        \cellcolor[rgb]{0.9,0.9,0.9}{
            \makecell[{{p{\linewidth}}}]{
                \texttt{\tiny{[GM$|$GM]}}
                \texttt{game successful} \\
            }
        }
    }
    & & \\ \\

    \theutterance \stepcounter{utterance}  
    & & & \multicolumn{2}{p{0.3\linewidth}}{
        \cellcolor[rgb]{0.9,0.9,0.9}{
            \makecell[{{p{\linewidth}}}]{
                \texttt{\tiny{[GM$|$GM]}}
                \texttt{end game} \\
            }
        }
    }
    & & \\ \\

\end{supertabular}
}

\end{document}
